\subsubsection{Separately Defined Stochastic Reference}
\label{stochastic-reference}

The principal controlled evaluation will examine receiving-state accuracy
under a specified stochastic process, alongside the index-behavior analysis.
This reference will adopt Independent Cascade's conditional edge-transmission
mechanism and live-edge interpretation \citep{kempe2003}. The study will add
independently sampled initial activations at all nodes. For this benchmark only,
the supplied numbers $B_j$ and $w_{ij}=S_{ij}\tau_{ij}$ will parameterize
independent Bernoulli variables:
\[
Z_j\sim\operatorname{Bernoulli}(B_j),\qquad
E_{ij}\sim\operatorname{Bernoulli}(w_{ij}),\qquad
X_j=Z_j\lor\bigvee_{i\in P(j)}(X_i\land E_{ij}).
\eqtag{3-10}
\]
All primitive $Z$ and $E$ variables will be mutually independent; derived
node activations need not be independent. The Boolean recursion will be
evaluated in topological order. A trial will draw each primitive once and
reuse its value wherever its consequences recur. The reference output will
be $\pi_j=\Pr(X_j=1)$, not expected spread obtained by selecting $j$ as a seed.

Conditional transmission will use $w_{ij}$, not $Q_{ij}=w_{ij}R_i$.
Sampling the upstream activation and then using $Q_{ij}$ would incorporate
the upstream marginal twice. GBBRPM will retain its deterministic index
definition. Differences $R_j-\pi_j$ will be called approximation errors only
within this explicitly probabilistic benchmark. Agreement will not calibrate
the domain adapters or establish real-world failure, exploitation, or
inspection utility.

For small fixtures with at most 16 non-degenerate primitive variables,
the exact solver will enumerate their binary configurations, evaluate the
Boolean recursion, and sum configuration probabilities for each activated
node. Probabilities of zero or one will be fixed, not enumerated. The solver
will check total probability mass and chain, isolated-node, independent-
convergence, and diamond identities before larger runs proceed. On graphs
whose incoming activation events have disjoint primitive ancestry, the
complement-product recursion will agree with this reference by independence
and induction. The study will state that sufficient condition explicitly;
structural overlap will not be assumed to make it necessary.

For larger graphs, the Monte Carlo solver will estimate
$\widehat\pi_j=N^{-1}\sum_{t=1}^{N}X_j^{(t)}$ from $N=100{,}000$
independent trials in memory-bounded batches. A predeclared seed will be
stored per graph and scenario. Sampling uncertainty will be reported using
simultaneous Hoeffding intervals \citep{hoeffding1963} $[\widehat\pi_j-h,\widehat\pi_j+h]$
clipped to $[0,1]$, where a two-sided concentration bound and a union bound over nodes give
\[
h=\sqrt{\frac{\log(2|V|/0.05)}{2N}}.
\eqtag{3-11}
\]
This conservative concentration bound applies to the sampling estimator
across nodes in one scenario; it is not an error bound for GBBRPM or a
simultaneous guarantee across every experiment. Monte Carlo will be checked
against enumeration on the eligible fixtures. Validation discrepancies and
interval coverage will be recorded, not silently removed. Graphs sharing
primitive draws will share the same reference for all deterministic methods.

\subsubsection{Benchmark Bracket and Its Interpretation}
\label{benchmark-bracket}

\textit{Proposition (receiving-state bracket under independent primitives).}
Let $G$ be a finite DAG and let all initial-activation and edge-transmission
variables in Equation~(3-10) be mutually independent, with supplied parameters
in $[0,1]$. On the same graph and inputs, define recursive maximum and
clipped-additive states by
\[
M_j=\max\left(B_j,\max_{i\in P(j)}w_{ij}M_i\right),\qquad
H_j=\min\left(1,B_j+\sum_{i\in P(j)}w_{ij}H_i\right),
\eqtag{3-12}
\]
where the incoming maximum is zero for an empty predecessor set. Let $R_j$
be ordinary GBBRPM and $\pi_j=\Pr(X_j=1)$ under the benchmark. Let $T_j$
be the exact activation probability on the target-specific restricted
in-arborescence defined in Section~\ref{reference-comparators}. It will
retain a consistent maximum-weight path from every positive-path ancestor,
the original local inputs and weights, and target $j$ itself; no path-weight
threshold will be applied. Then
\[
M_j\le T_j\le\pi_j\le R_j\le H_j\qquad\text{for every }j\in V.
\eqtag{3-13}
\]

\textit{Proof.} In a DAG, $X_i$ cannot depend on the primitive $E_{ij}$,
and no incoming activation at $j$ can depend on $Z_j$. Independence of
the primitives therefore gives
$\Pr(X_i\land E_{ij})=w_{ij}\pi_i$. Since activation at $j$ includes
its own initial event and every successful incoming event,
$\pi_j\ge\max(B_j,\max_i w_{ij}\pi_i)$. A predecessor-free node has
$M_j=\pi_j=B_j$; induction over topological order establishes the lower bound.

The maximum recursion also has the path representation
$M_j=\max_{a,\gamma:a\rightarrow j}B_a\prod_{e\in\gamma}w_e$,
including the zero-length path at $j$. The restricted tree retains a
maximum-weight path for each positive-path ancestor by consistent next-edge
choices. Activation at $a$ through its own local event followed by successful
transmission along that retained path is sufficient to activate $j$ in the
tree. Its probability is $B_a\prod_{e\in\gamma}w_e$ by primitive independence.
The tree probability is at least each such path contribution and $B_j$, so
$M_j\le T_j$. Ancestors without a positive path contribute zero to $M_j$.
Couple the restricted and full processes using the same local and edge draws.
Deleting edges cannot create an activation, hence $T_j\le\pi_j$.
On the target-directed tree, incoming ancestor subtrees have disjoint
primitives; complement-product recursion therefore evaluates $T_j$ exactly.
This argument relies on the stated unthresholded, path-preserving construction;
arbitrary pruning need not retain the $T_j\ge M_j$ guarantee.

For the upper bound, let $F_{ij}$ denote failure of incoming activation
through $(i,j)$, that is, the complement of $\{X_i=1,E_{ij}=1\}$.
Every $X_i$ is an increasing Boolean function of the primitive variables,
so every $F_{ij}$ is decreasing. Positive association under a product
Bernoulli measure, the Harris/FKG principle, gives
$\Pr(\bigcap_iF_{ij})\ge\prod_i\Pr(F_{ij})$. This established principle
is also used for cascade upper bounds by \citet{abbe2017,abbe2020}; it is not a
new correlation inequality. Since $Z_j$ is independent of the incoming
events,
\[
\pi_j=1-(1-B_j)\Pr\left(\bigcap_{i\in P(j)}F_{ij}\right)
\le1-(1-B_j)\prod_{i\in P(j)}(1-w_{ij}\pi_i).
\eqtag{3-14}
\]
Complement product is increasing in each predecessor marginal.
Starting from $\pi_j=R_j=B_j$ at predecessor-free nodes, topological
induction yields $\pi_j\le R_j$.

Finally, for nonnegative inputs in $[0,1]$,
$1-\prod_k(1-x_k)\le\min(1,\sum_kx_k)$. Applying this inequality to
the local term and incoming contributions, followed by topological
induction and monotonicity of clipped addition, gives $R_j\le H_j$.
Primitive probabilities of zero or one follow directly or by continuity.
This completes the argument.

The proposition will be presented as a specialization of established
association and bounded-aggregation arguments, not as a verified first
theorem about cascade approximation. Abbe et al.'s journal treatment
also bounds individual activation probabilities before summing them into
seed-set influence. Its upper method uses length-indexed messages; its
DAG-based lower method uses a different incoming-message recurrence
\citep{abbe2020}. Section~\ref{cascade-bounds-prior-art} compares the
constructions and limits. Independent local activation
can also be represented by one certainly active auxiliary source with
independent edges of probabilities $B_j$ to all original nodes. That
construction facilitates comparison with fixed-seed literature; it will
not conceal the model's mathematical precedents.

For this benchmark, score discrepancy is one-sided: $R_j-\pi_j\ge0$.
Its scalar certificate is $0\le R_j-\pi_j\le R_j-M_j$.
When the tree comparator is computed, the tighter certificate is
$0\le R_j-\pi_j\le R_j-T_j\le R_j-M_j$.
Computing $T_j$ requires tree construction and evaluation for each requested
target; the tighter certificate is not free relative to the scalar pass.
Its reporting will reuse the already computed comparator result.
GBBRPM is no looser than recursive clipped addition as a nodewise upper
bound; this ordering need not be rediscovered experimentally. Experiments
will measure bound width, actual slack, conditions of equality, and
ranking or selection effects. A smaller pointwise score error does not
guarantee better rank order. The condition $T_a>R_b$ certifies
$\pi_a>\pi_b$ when tree bounds are available; $M_a>R_b$ remains a cheaper
sufficient certificate. Overlapping brackets remain inconclusive,
not evidence of a reversed ordering.

These guarantees concern only the stated independent-primitive benchmark.
They do not certify probabilities or physical bounds for the domain index,
dependent initial events, correlated transmission primitives, feedback, or
other cascade mechanisms. The ancestry-aware diagnostic will require its
own analysis and will not inherit the GBBRPM upper-bound guarantee.


\subsubsection{Matched Dependence Controls and Graph Inventory}
\label{dependence-controls}

The first experiment will separate uncertain common origins from independent
origins on a fixed graph. A root $a$ will connect with weight one to $m$
branch-start nodes $s_1,\ldots,s_m$. Each branch start will connect to a
common receiver $d$ through an edge-disjoint path of $\ell$ edges, all with
weight $p$. Intermediate local disturbances and $B_d$ will be zero. In the
shared-origin condition, $B_a=b$ and all branch-start local terms will be
zero. In the independent-origin condition, $B_a=0$ and every $B_{s_r}=b$.
The graph, weights, and branch-start activation marginals will match; only
the joint source construction will change. Both conditions will include
the inactive root in their node inventory.

For $u=p^\ell$, this constructed family has exact receiver references
\[
\pi_d^{\mathrm{shared}}=b\{1-(1-u)^m\},\qquad
\pi_d^{\mathrm{independent}}=1-(1-bu)^m.
\eqtag{3-15}
\]
GBBRPM will give $1-(1-bu)^m$ in both conditions. This identity will provide
an analytical benchmark, not a claim of a newly invented dependence theory.
For two branches, the shared-origin discrepancy is $b(1-b)u^2$.
Thus shared ancestry alone does not force discrepancy: $b=1$ gives
agreement in this family. With $b=0.5$ and $u=1$, the reference is $0.5$
and GBBRPM is $0.75$. With $b=1$ and $u=0.25$, both are $0.4375$;
retaining only the strongest branch gives $0.25$. These algebraic controls
will prevent strongest-path pruning from being assumed correct in advance.

\begin{table}[H]
\centering
\caption{Planned reference-benchmark inventory, separate from historical runs.}
\label{tab:reference-inventory}
\small
\begin{tabularx}{\textwidth}{@{}L L r@{}}
\toprule
Family & Fixed protocol & Scenarios\\
\midrule
Matched branch family & Four branch counts, two path lengths, seven
edge weights, five source probabilities, and two origin conditions & 560\\
Generated DAGs & Three families, four sizes, 25 graph seeds, six weights,
and two initial-activation settings & 3600\\
Express topology & One verified snapshot, six weights, and two controlled
initial-activation settings & 12\\
\bottomrule
\end{tabularx}
\end{table}

The branch-family grid will use $m\in\{1,2,4,8\}$,
$\ell\in\{2,4\}$, $p\in\{0.05,0.10,0.20,0.25,0.50,0.75,1\}$,
and $b\in\{0.10,0.25,0.50,0.75,1\}$, under both origin conditions.
The closed form will supply exact receiver references at every size,
including rare-event regimes; the 16-variable enumeration cap will remain
an implementation-validation choice rather than the limit of exact evidence.
Conditioning or decision-diagram solvers will be optional future tooling,
not required to establish these references.

The generated-DAG stage will use $|V|\in\{20,50,100,250\}$ and three
explicit constructions. Each will have $2|V|$ unique edges and a recorded
topological ordering:
\begin{enumerate}[label=\alph*.,leftmargin=*]
\item A forward-edge family will sample edges uniformly from all ordered
pairs consistent with a seeded random node ordering.
\item A layered family will distribute nodes as evenly as possible over
$\lceil\sqrt{|V|}\rceil$ ordered layers, randomize order within each layer,
and sample edges uniformly between adjacent layers. The available-pair
count will be checked before sampling; no hidden generator fallback will
replace an infeasible configuration.
\item A preferential-attachment DAG will start with all ten forward edges
among five ordered nodes. Each later node will choose two distinct earlier
predecessors without replacement, proportional to their current total degree
plus one, and receive edges from them. This supplies exactly $2|V|$ edges.
The construction will be called preferential attachment; a scale-free degree
law will not be asserted without evidence.
\end{enumerate}

Per family and size, seeds 41--45 will supply five development graphs and
1001--1020 will supply 20 evaluation graphs, yielding 300 distinct topologies.
Each topology will retain its edges under
$w\in\{0.05,0.10,0.20,0.25,0.50,0.75\}$. The first
$\lceil0.10|V|\rceil$ nodes in its ordering will receive either $B=0.5$
or $B=1$; other local terms will be zero. These changes will create 3600
scenarios, of which 720 are development and 2880 are evaluation scenarios.
A local-origin node need not have in-degree zero. The controlled uniform
weights will isolate topology, dependence, and attenuation; heterogeneous
N1--N5 inputs will remain a separate information-construction comparison.

The record will include depth, degree distributions, reachability, source
coverage, and the fraction of receivers with overlapping incoming local-origin
lineages. Structural overlap will not be equated with statistical dependence
strength. Family comparisons will be descriptive, because depth and degree
patterns differ even with equal node and edge counts. The fixed branch
family will supply the controlled dependence manipulation.

A further 12 scenarios will use the Express 4.18.2 topology, subject to
recovering a verified snapshot. The historical manifest records hashes for
\texttt{nodes.csv} and \texttt{edges.csv}, but those raw inputs are absent
from this manuscript package. The original files will be recovered and hash-
checked, or a new resolved-dependency snapshot will be acquired, versioned,
and explicitly distinguished from the historical 71-node, 128-edge case.
Edges will run from dependency to dependent; acyclicity and the reported
shared paths will be checked. Cycles will not be silently deleted.

For each of the six weights above, the seven historical advisory-origin
identities will receive controlled initial probabilities of either $0.5$
or $1$, with other local terms zero. Their identities will be verified from
the source records; real advisory severity values will not be interpreted
as calibrated probabilities. This will be a stochastic experiment on real
topology, not incident validation. If required records cannot be recovered,
the case will be reported unavailable rather than reconstructed from figures.
It will remain separate from generated-graph averages.

At $N=100{,}000$ trials, small receiving probabilities may be unresolved.
The study will therefore preselect 48 precision-audit scenarios: evaluation
seed 1001 for every family and size, $w\in\{0.05,0.10\}$, and both
initial-activation settings. Each will receive an additional independent
$N=1{,}000{,}000$-trial reference with its own interval. Both estimates
will be retained; this subset will be selected before outcomes are inspected.
Coverage statements will remain per solve, not a joint guarantee for all
solves. Zero observed activations will not be called zero true probability.
Elsewhere unresolved differences will remain unresolved. Batching will cap
memory; development timing will determine whether the full prospective
inventory is feasible. Any reduction will be documented and frozen before
evaluation outcomes, including its resulting scenario counts.

\subsubsection{Task-Matched Comparators and Recursion Ablation}
\label{reference-comparators}

Every deterministic method will receive the same graph, $B$, and effective
weights $w$. The comparison will include local-only scoring and the three
recursive aggregators in Equation~(3-9). A recursion ablation will retain
the complement-product operator but substitute original local input for
current upstream state:
\[
R_j^{\mathrm{local\text{-}upstream}}
=1-(1-B_j)\prod_{i\in P(j)}(1-w_{ij}B_i).
\eqtag{3-16}
\]
This condition will isolate inherited multihop state rather than be called
an equivalent cascade solver. Recursive maximum aggregation will supply
the strongest local-origin weighted-path comparator; an additional
strongest-path method would duplicate it under these inputs. An arithmetic
mean will not be required because adding a zero input can reduce its state,
contrary to the selected accumulation requirement.

A MIA-inspired restricted-in-arborescence comparator will also be evaluated
\citep{chen2010mia}. For each target $j$, reverse topological dynamic
programming will compute the maximum positive weighted-path product from
every ancestor to $j$. Each ancestor will retain one next edge on a maximizing
path, with exact ties broken consistently by node identifier. The union of
these target-directed paths will form an in-arborescence. No positive-path
threshold will be applied in the primary comparator. Original local $B$
terms on retained nodes and the retained edge weights will be preserved.
Complement-product recursion on that tree will compute the target's exact
activation probability for the restricted independent-primitive process.
Edge deletion makes this a lower bound on the full benchmark marginal.
This will be an adaptation of MIA's tree-restriction idea to fixed receiving
queries, not its complete greedy seed-selection algorithm. Building one tree
per target will have a different cost from a single global scalar pass.

The ancestry-aware diagnostic in Section~\ref{behavioral-evaluation} will
also be compared with the same reference. Each recursive variant will
compute its own upstream states. In-degree will remain a descriptive
structural comparison; descendant count and seed-spread measures will not
be used as receiving-probability accuracy baselines without a separately
specified task. Exact and Monte Carlo methods will be reference solvers,
with their precision and computational cost stated.

\subsubsection{Bound Tightness and Fixed-Budget Ranking Metrics}
\label{reference-metrics}

Fixed-budget selection and ranking will be the primary benchmark outcomes;
score evaluation will characterize the analytic bracket's tightness.
Per scenario, the study will report width $R_j-M_j$, actual upper slack
$R_j-\pi_j$, lower slack $\pi_j-M_j$, mean and maximum absolute error,
and terminal error relative to exact or estimated marginals. The guaranteed
nonnegative GBBRPM slack will not be treated as an untested hypothesis.
Numerical bound violations against exact references will be investigated
as implementation or assumption defects. Node averages will first be computed within a graph;
cross-graph summaries will give graphs equal weight and show their spread.
Graph families, sizes, weight regimes, and source settings will be reported
separately. Each graph will remain the replication unit; its 12 weight/source
conditions will not be treated as independent graphs. Evaluation summaries
will include paired method differences and their distribution across the
20 graph replicates. Development summaries will remain separate.
Monte Carlo intervals will accompany score discrepancies. For example,
a Monte Carlo interval $[L_j,U_j]$ can be intersected with the analytic
bracket $[M_j,R_j]$. If nonempty, $[L'_j,U'_j]$ gives possible upper
slack $[R_j-U'_j,R_j-L'_j]$. An empty intersection will be recorded as a
sampling-coverage or implementation discrepancy, not silently repaired. A difference smaller than reference uncertainty
will not be interpreted as a resolved advantage. Reports will include the
fractions of estimated reference states and model scores at or above $0.95$,
the existing numerical-saturation fraction, and the number of distinct score
groups. Near-saturation and near-zero cases will not be discarded because
they give inconvenient rankings. Bracket-certified pairwise orderings will
be counted separately from Monte Carlo-based rank estimates.

Ranking evaluation will exclude nodes with positive local inputs, so that
high supplied disturbances do not trivially dominate a receiving-state task.
The common candidate set $C$ will contain all remaining nodes, including
unreachable nodes. The study will report tie-aware Spearman association,
and selection at budgets $k=\max(1,\lceil f|C|\rceil)$ for
$f\in\{0.10,0.20,0.30\}$. Exactly $k$ nodes will be selected. Boundary
ties will be handled by averaging over 100 seeded random node-order
permutations, shared by all methods and unrelated to reference outcomes. Permutations will
use seed 140 and the $10^{-12}$ anchored score-tie policy stated in
Section~\ref{analysis-reproducibility}.
This policy will differ from the expanded cutoff sets used for descriptive
Jaccard comparisons in Section~\ref{analysis-reproducibility}.

For a selected set $T_k$, the benchmark captured activation mass and
selection regret will be
\[
\begin{aligned}
\operatorname{Capture}(T_k)&=
\frac{\sum_{j\in T_k}\pi_j}{\sum_{j\in C}\pi_j},\\
\operatorname{Regret}_k&=
\sum_{j\in T_k^*}\pi_j-\sum_{j\in T_k}\pi_j,
\end{aligned}
\eqtag{3-17}
\]
where $T_k^*$ maximizes activation mass over $k$ candidates. Capture will
be undefined when the denominator is zero; the all-zero case will be
reported explicitly. Monte Carlo-based oracle membership and regret will
be labeled estimates. A conservative bound for the difference in mass
between two selected sets is $h$ times their symmetric-difference size;
results unresolved at that precision will remain unresolved. These metrics
concern expected node activation, not avoided loss or intervention benefit.

The main comparisons will be paired by graph and scenario. They will ask
when the analytic bounds are tight and whether complement-product recursion
improves receiving-state selection over simpler rules or the restricted-tree
comparator. Clipped sum is analytically a looser upper bound; ranking superiority
does not follow from that fact. The study will also ask whether the dependence diagnostic
improves either measure without unacceptable resource growth. Negative,
equivalent, and mixed findings will be retained. The study will describe
an empirical applicability boundary rather than presume a correction wins
or claim universal probabilistic accuracy for the index.
